\section{Conclusion}

This work presents a calibration-free pipeline that estimates the root canal filling length directly from preoperative periapical radiographs. The method combines single-tooth extraction, image enhancement, lightweight instance segmentation, and small-sample regression to address the missing pixel scale, inconsistent export scales, and low apical contrast of clinically exported images. On an independent test set of 67 anterior teeth, the corrected MAE was 0.993~mm with an ideal ratio of 62.7\% and a short overall bias of $-0.493$~mm, outperforming CLAHE and unprocessed images.

The pipeline first extracts the target tooth from the full radiograph with a YOLO11n detector and standardises the input size with aspect-ratio-preserving scaling and padding. A morphological top-hat transformation then enhances local contrast, and YOLO11n-seg produces the tooth mask, along whose principal axis the pixel length is measured. Finally, long-side normalisation removes scale differences, and GPR with patient age converts the length to millimetres before a 0.5-mm safety offset is subtracted.

Overall, the proposed pipeline balances estimation accuracy and clinical safety on the evaluated dataset and provides a practical technical basis for preoperative root-length assessment without a physical scale. Future work will include more tooth types and samples, multicentre data, and the use of acquisition parameters or lightweight calibration information to further improve accuracy, together with prospective clinical studies to confirm its value for preoperative assessment.
